\chapter{Introduction}
\label{chap:introduction}

Suppose two independent groups are observed on the same categorical variables
\(X\) and \(Y\). Mutual information (MI) describes the strength of association
between those variables in each group. If the estimated MI differs, the
difference alone does not tell us whether the populations differ: both estimates
vary when new observations are drawn. The question in this thesis is whether
the two \emph{population} MI values can plausibly be equal.

This is not the usual test of whether \(X\) and \(Y\) are independent within
one group. The null hypothesis here permits dependence in both groups and
permits their complete joint distributions to differ. It specifies only
\begin{equation}
  H_0:I(P)=I(Q),
  \qquad H_1:I(P)\ne I(Q),
  \label{eq:introduction-null}
\end{equation}
where \(P\) and \(Q\) are the two joint distributions. The samples from
\(P\) and \(Q\) are independent of each other; that sampling assumption does
not imply independence of \(X\) and \(Y\) within either population.

\section{Why the comparison is difficult}

MI is a nonlinear function of all cell probabilities and their row and column
margins. In a small or uneven table, many cells can be empty in the sample
even though their population probabilities are positive. The same data are
then used to estimate both the MI difference and its standard error. A
first-order normal Wald test is simple, but its finite-sample rejection rate
need not be close to the nominal level.

Welch--Satterthwaite inference offers a possible adjustment. Instead of
comparing the standardised difference with a normal reference, it compares
the same statistic with a Student reference whose effective degrees of
freedom describe uncertainty in the estimated variance. For MI, that variance
is itself a nonlinear function of the full table. A conventional \(n-1\)
assignment does not account for the way a changed cell also changes the
margins and the pointwise MI values. The method examined here estimates those
degrees of freedom from the sensitivity of the complete variance functional.

The resulting reference is an approximation, not a guarantee of better
inference. A heavier-tailed reference can reduce excessive false positives,
but it can also suppress useful detection when Wald already rejects too
rarely. Moreover, one method can fail to produce a valid p-value on samples
where the other succeeds. Calibration, detection and valid-result frequency
must therefore be considered together.

\section{Research questions and contributions}

The thesis asks four related questions. First, how can the sampling
uncertainty of the estimated MI variance be turned into an MI-specific
Welch--Satterthwaite reference? Second, how does that reference affect false
positives, detection and numerical validity relative to Normal Wald? Third,
how do table size, margins, baseline dependence, the location of association
and unequal sample sizes change the comparison? Fourth, do the observed
problems diminish with larger samples, and what additional computation does
the correction require?

The contribution is a connected derivation and a controlled evaluation of its
limits. The mathematical ingredients of first-order MI inference and
Satterthwaite moment matching are established; the thesis does not claim to
invent either one. It derives the sensitivity of the complete MI variance,
implements the resulting two-sample procedure, and evaluates both methods on
fixed populations whose true MI values are known. The experiment was frozen
after exploratory work and covers \(2\times2\) through \(8\times8\) square
tables, two rectangular shapes and explicit stress regimes.

The broad result is conditional rather than a uniform improvement. Expanded
Welch corrects some liberal null rejection in small samples. It can also be
too conservative, especially where Wald is already below the 5\% target, and
can produce fewer valid results in sparse samples. The final evidence does
not support replacing Wald generally. This negative finding is useful: it
identifies what the proposed adjustment changes, what it cannot change, and
which conditions need more than a new reference quantile.

\section{Scope and organisation}

The analysis concerns independent multinomial samples on fixed, aligned
categorical alphabets. The principal theory assumes strictly positive cell
probabilities and positive first-order MI variances. Exact independence is
studied separately as a boundary case, not treated as an ordinary positive-MI
regime. The final experiment compares Normal Wald with Expanded Welch at a
two-sided nominal level of 0.05. It does not establish performance for paired
data, continuous MI estimators or alphabets that grow with sample size.

Chapter~\ref{chap:background} supplies the statistical background.
Chapters~\ref{chap:baselines} and~\ref{chap:expanded} define the shared
statistic and derive the correction. Chapter~\ref{chap:design} explains how
the populations and samples are generated and how the results are measured.
Chapter~\ref{chap:results} shows the complete landscape before the focused
comparisons. Chapters~\ref{chap:discussion} and~\ref{chap:conclusion} discuss
the consequences and the resulting recommendation. The appendices provide
the detailed algebra, construction and source-to-result record.
